\documentclass[a4paper,12pt]{article}

\usepackage[T2A]{fontenc}
\usepackage[utf8]{inputenc}
\usepackage[russian]{babel}
\usepackage{amsmath,amssymb,amsthm,mathtools}
\usepackage{helvet}
\renewcommand{\familydefault}{\sfdefault}
\usepackage{icomma}
\usepackage[a4paper,margin=2.15cm,headheight=15pt]{geometry}
\usepackage{microtype}
\usepackage{tikz}
\usepackage{tkz-euclide}
\usepackage{pgfplots}
\pgfplotsset{compat=1.18}
\usepackage{tabularx,booktabs,longtable}
\usepackage{enumitem}
\usepackage{hyperref}
\usepackage{parskip}
\usepackage{fancyhdr}
\usepackage{setspace}
\usepackage{xcolor}

% -------------------------------------------------
% Цветовая палитра
% -------------------------------------------------
\definecolor{accent}{HTML}{008080}
\definecolor{darkaccent}{HTML}{004D4D}
\definecolor{textgray}{HTML}{333333}
\definecolor{softgray}{HTML}{EEF3F3}
\definecolor{linegray}{HTML}{B7CACA}
\definecolor{warn}{HTML}{8A4D3B}

% -------------------------------------------------
% Общая типографика
% -------------------------------------------------
\onehalfspacing
\setlength{\parindent}{0pt}
\setlength{\parskip}{0.55em}
\setlist[itemize]{leftmargin=1.6em,itemsep=0.25em,topsep=0.25em}
\setlist[enumerate]{leftmargin=1.9em,itemsep=0.55em,topsep=0.35em}

\hypersetup{
    colorlinks=true,
    linkcolor=darkaccent,
    urlcolor=darkaccent,
    pdftitle={Алкены: строение, изомерия и номенклатура},
    pdfauthor={Лёвин Артём Александрович}
}

\pagestyle{fancy}
\fancyhf{}
\fancyhead[L]{\small\textcolor{textgray}{Алкены: строение, изомерия и номенклатура}}
\fancyhead[R]{\small\textcolor{textgray}{Ксения Клыкова}}
\fancyfoot[C]{\small\textcolor{textgray}{\thepage}}
\renewcommand{\headrulewidth}{0.4pt}
\renewcommand{\headrule}{\hbox to\headwidth{\color{linegray}\leaders\hrule height \headrulewidth\hfill}}

% -------------------------------------------------
% Вспомогательные команды
% -------------------------------------------------
\newcommand{\lessonsection}[1]{%
    \vspace{0.7em}
    {\Large\bfseries\color{darkaccent}#1}\par
    \vspace{0.15em}
    {\color{linegray}\rule{\linewidth}{0.7pt}}
    \vspace{0.35em}
}

\newcommand{\sublesson}[1]{%
    \vspace{0.4em}
    {\large\bfseries\color{accent}#1}\par
}

\newcommand{\important}[1]{\textbf{\color{darkaccent}#1}}
\newcommand{\warning}[1]{\textbf{\color{warn}#1}}

\begin{document}

% =================================================
% ТИТУЛЬНЫЙ ЛИСТ
% =================================================
\begin{titlepage}
\thispagestyle{empty}

\begin{center}
\vspace*{2.5cm}

{\Large\color{accent}\textbf{ЧЕК-ЛИСТ}}

\vspace{1.1cm}

{\fontsize{25}{31}\selectfont\bfseries\color{darkaccent}
Алкены}

\vspace{0.45cm}

{\Large\bfseries\color{textgray}
Строение, изомерия\\[0.2em]
и правила номенклатуры}

\vspace{1.8cm}

{\large
Ключевые идеи первого занятия по теме}

\vspace{2.2cm}

{\large\bfseries
Лёвин Артём Александрович\\
эксклюзивно для Ксении Клыковой}

\vspace{1.1cm}

{\large 10 класс}

\vfill

{\large \today}

\vspace{1.6cm}

\begin{tikzpicture}[x=1cm,y=1cm]
    \draw[
        color=accent,
        line width=1.15pt
    ]
    (0,0)
    .. controls (3.0,0.65) and (7.0,-0.65) ..
    (10,0)
    .. controls (12.2,0.35) and (14.0,0.25) ..
    (15.6,0);
\end{tikzpicture}

\end{center}
\end{titlepage}

% =================================================
% ОСНОВНОЙ ЧЕК-ЛИСТ
% =================================================

\lessonsection{1. Что такое алкены}

\important{Алкены} --- ациклические непредельные углеводороды, содержащие одну двойную связь
\[
\mathrm{C=C}.
\]

Для алкенов с одной двойной связью и без циклов справедлива общая формула
\[
\boxed{\mathrm{C_nH_{2n}},\qquad n\ge 2.}
\]

Например:
\[
\mathrm{C_2H_4,\quad C_3H_6,\quad C_4H_8,\quad C_5H_{10}.}
\]

Первый представитель гомологического ряда:
\[
\mathrm{CH_2=CH_2}
\qquad\text{--- этен (этилен).}
\]

\sublesson{Что важно сразу проверять}

Если перед Вами два предполагаемых изомера, сначала сравните их \important{молекулярные формулы}.

Например,
\[
\mathrm{C_5H_{10}}
\qquad\text{и}\qquad
\mathrm{C_5H_9}
\]
не могут быть изомерами, потому что их состав различен.

\[
\boxed{
\text{Изомеры имеют одинаковый качественный и количественный состав,}
}
\]
\[
\boxed{
\text{но различное строение молекул.}
}
\]

% -------------------------------------------------

\lessonsection{2. Изомерия алкенов}

Для изученных на занятии алкенов особенно важны два вида структурной изомерии.

\sublesson{2.1. Изомерия углеродного скелета}

Молекулярная формула сохраняется, но меняется строение углеродной цепи.

Например:
\[
\mathrm{CH_2=CH-CH_2-CH_2-CH_3}
\]
и
\[
\mathrm{CH_2=C(CH_3)-CH_2-CH_3}.
\]

Оба вещества имеют формулу
\[
\mathrm{C_5H_{10}},
\]
но углеродный скелет устроен по-разному.

\sublesson{2.2. Изомерия положения двойной связи}

Углеродный скелет может оставаться тем же, но положение связи \(\mathrm{C=C}\) меняется.

Например:
\[
\mathrm{CH_2=CH-CH_2-CH_2-CH_3}
\]
--- пентен-1,

\[
\mathrm{CH_3-CH=CH-CH_2-CH_3}
\]
--- пентен-2.

Оба вещества имеют формулу
\[
\mathrm{C_5H_{10}},
\]
однако двойная связь занимает разное положение.

\sublesson{Критерий изомерии}

Чтобы два изображения действительно соответствовали разным изомерам, должны выполняться одновременно два условия:

\begin{enumerate}
    \item молекулярная формула одинакова;
    \item порядок соединения атомов различен.
\end{enumerate}

\warning{Важно:} два рисунка одной и той же молекулы, полученные её поворотом или разворотом на плоскости, не являются разными изомерами.

% -------------------------------------------------

\lessonsection{3. Как выбрать главную цепь}

При составлении названия алкена используйте следующий порядок.

\begin{enumerate}
    \item Найдите углеродную цепь, \important{обязательно содержащую двойную связь}.
    \item Среди таких цепей выберите наиболее длинную.
    \item Пронумеруйте её с того конца, который ближе к двойной связи.
    \item Определите положение двойной связи.
    \item Найдите все заместители и их положения.
    \item Составьте полное название вещества.
\end{enumerate}

\sublesson{Принципиально важное правило}

\[
\boxed{
\text{Цепь без двойной связи не может быть главной цепью алкена,}
}
\]
даже если она содержит больше атомов углерода.

Иными словами, наличие связи \(\mathrm{C=C}\) имеет приоритет при выборе основной цепи.

% -------------------------------------------------

\lessonsection{4. Нумерация главной цепи}

Главную цепь нумеруют с того конца, от которого двойная связь получает
\important{наименьший возможный номер}.

Например:
\[
\mathrm{CH_3-CH_2-CH=CH_2}.
\]

Правильная нумерация начинается справа:
\[
\mathrm{CH_3-CH_2-CH=CH_2}
\]
\[
\hspace{0.2cm}4\qquad 3\qquad 2\qquad 1
\]

Поэтому вещество называется
\[
\boxed{\text{бутен-1}},
\]
а не бутен-3.

\sublesson{Почему направление рисунка не определяет название}

Структурную формулу можно записать в противоположном направлении. Например,
\[
\mathrm{CH_2=C(CH_3)-CH_3}
\]
и
\[
\mathrm{CH_3-C(CH_3)=CH_2}
\]
изображают одно и то же вещество.

В обоих случаях после правильного выбора главной цепи и её нумерации получается
\[
\boxed{\text{2-метилпропен-1}.}
\]

Следовательно, само расположение формулы слева направо не определяет строение вещества: решающим является порядок соединения атомов.

% -------------------------------------------------

\lessonsection{5. Как образуются названия алкенов}

Название основной цепи определяется числом атомов углерода.

\begin{center}
\renewcommand{\arraystretch}{1.3}
\begin{tabular}{ccc}
\toprule
Число атомов C & Корень & Пример алкена \\
\midrule
2 & эт-   & этен \\
3 & проп- & пропен \\
4 & бут-  & бутен \\
5 & пент- & пентен \\
6 & гекс- & гексен \\
\bottomrule
\end{tabular}
\end{center}

Наличие одной двойной связи отражается суффиксом
\[
\boxed{\text{-ен}.}
\]

Положение двойной связи указывают номером первого атома углерода этой связи.

Например:
\[
\mathrm{CH_2=CH-CH_2-CH_2-CH_3}
\]
--- \important{пентен-1},

\[
\mathrm{CH_3-CH=CH-CH_2-CH_3}
\]
--- \important{пентен-2}.

В школьной номенклатуре часто используется запись
\[
\text{пентен-1},\qquad \text{пентен-2}.
\]

% -------------------------------------------------

\lessonsection{6. Заместители}

Углеводородные ответвления в главной цепи рассматривают как
\important{алкильные заместители}.

Основные заместители, встречавшиеся на занятии:

\begin{center}
\renewcommand{\arraystretch}{1.35}
\begin{tabularx}{0.82\linewidth}{l c X}
\toprule
Заместитель & Формула & Название \\
\midrule
метильный & \(\mathrm{-CH_3}\) & метил \\
этильный & \(\mathrm{-CH_2-CH_3}\) & этил \\
\bottomrule
\end{tabularx}
\end{center}

Перед названием заместителя обязательно указывают номер атома главной цепи, при котором он расположен.

Например:
\[
\mathrm{CH_2=C(CH_3)-CH_3}
\]
--- \important{2-метилпропен-1}.

Здесь:
\begin{itemize}
    \item главная цепь содержит \(3\) атома углерода;
    \item двойная связь начинается у атома \(1\);
    \item метильная группа находится у атома \(2\).
\end{itemize}

% -------------------------------------------------

\lessonsection{7. Несколько одинаковых заместителей}

Если одинаковых заместителей несколько, используют приставки:

\begin{center}
\renewcommand{\arraystretch}{1.25}
\begin{tabular}{cc}
\toprule
Количество & Приставка \\
\midrule
2 & ди- \\
3 & три- \\
4 & тетра- \\
\bottomrule
\end{tabular}
\end{center}

При этом необходимо указать положение \important{каждого} заместителя.

Например, если метильные группы расположены при атомах \(2\), \(3\) и \(4\), фрагмент названия записывают как
\[
\boxed{\text{2,3,4-триметил-}\ldots}
\]

Нельзя записывать только один локант: все положения заместителей должны быть отражены в названии.

% -------------------------------------------------

\lessonsection{8. Разные заместители}

Если в молекуле присутствуют разные заместители, их названия перечисляют
\important{в алфавитном порядке}.

Например, среди заместителей
\[
\text{метил},\qquad \text{этил}
\]
сначала в названии указывают \important{метил}, затем \important{этил}, поскольку русская буква «м» стоит в алфавите раньше буквы «э».

При алфавитном сравнении приставки количества
\[
\text{ди-, три-, тетра-}
\]
не учитывают.

% -------------------------------------------------

\lessonsection{9. Универсальный алгоритм названия алкена}

\begin{center}
\begin{tikzpicture}[
    node distance=7mm,
    every node/.style={
        align=center,
        text width=11.2cm,
        minimum height=8mm,
        font=\small
    },
    step/.style={
        draw=accent,
        rounded corners=2pt,
        line width=0.7pt,
        inner sep=5pt
    },
    arrow/.style={
        ->,
        >=stealth,
        color=accent,
        line width=0.8pt
    }
]

\node[step] (a) {1. Найдите цепь, содержащую двойную связь};
\node[step, below of=a] (b) {2. Выберите наиболее длинную из подходящих цепей};
\node[step, below of=b] (c) {3. Пронумеруйте её от конца, ближайшего к \(\mathrm{C=C}\)};
\node[step, below of=c] (d) {4. Укажите положение двойной связи};
\node[step, below of=d] (e) {5. Найдите заместители и определите их локанты};
\node[step, below of=e] (f) {6. Для одинаковых заместителей используйте ди-, три-, тетра-};
\node[step, below of=f] (g) {7. Разные заместители перечислите по алфавиту};
\node[step, below of=g] (h) {8. Проверьте название повторной нумерацией структуры};

\draw[arrow] (a) -- (b);
\draw[arrow] (b) -- (c);
\draw[arrow] (c) -- (d);
\draw[arrow] (d) -- (e);
\draw[arrow] (e) -- (f);
\draw[arrow] (f) -- (g);
\draw[arrow] (g) -- (h);

\end{tikzpicture}
\end{center}

% -------------------------------------------------

\lessonsection{10. Обратная задача: от названия к формуле}

Если дано название вещества, действуйте в обратном порядке.

Например:
\[
\text{2-метилпентен-1}.
\]

\begin{enumerate}
    \item Корень \textbf{пент-} означает главную цепь из пяти атомов углерода.
    \item Окончание \textbf{-ен} указывает на двойную связь.
    \item Число \(1\) означает:
    \[
    \mathrm{C_1=C_2}.
    \]
    \item Запись \textbf{2-метил-} означает заместитель \(\mathrm{-CH_3}\) при втором атоме основной цепи.
\end{enumerate}

Получаем:
\[
\boxed{\mathrm{CH_2=C(CH_3)-CH_2-CH_2-CH_3}.}
\]

Проверка числа атомов:
\[
\mathrm{C_6H_{12}},
\]
что соответствует общей формуле алкенов
\[
\mathrm{C_nH_{2n}}
\]
при \(n=6\).

% -------------------------------------------------

\lessonsection{11. Проверка построенной структуры}

После построения структурной формулы полезно выполнить четыре быстрые проверки.

\begin{enumerate}
    \item \important{Главная цепь.} Число её атомов соответствует корню названия?
    \item \important{Двойная связь.} Она действительно начинается у указанного атома?
    \item \important{Заместители.} Каждый находится при требуемом атоме?
    \item \important{Валентность.} У каждого атома углерода суммарно четыре связи?
\end{enumerate}

Для ациклического моноалкена дополнительно проверьте молекулярную формулу:
\[
\boxed{\mathrm{C_nH_{2n}}.}
\]

% -------------------------------------------------

\lessonsection{12. Типичные ошибки}

\begin{enumerate}
    \item \warning{Выбирать самую длинную цепь вообще, не проверяя наличие двойной связи.}

    Правильно: главная цепь алкена обязана содержать \(\mathrm{C=C}\).

    \item \warning{Нумеровать с ближайшего к заместителю конца.}

    Правильно: в рассмотренных алкенах приоритет при выборе направления нумерации имеет положение двойной связи.

    \item \warning{Нумеровать атомы углерода заместителя вместе с главной цепью.}

    Правильно: номера получают только атомы выбранной главной цепи.

    \item \warning{Считать зеркально или развёрнуто записанную формулу новым веществом.}

    Правильно: сравните порядок соединения атомов и систематическое название.

    \item \warning{Забывать локанты при нескольких одинаковых заместителях.}

    Нужно писать, например,
    \[
    \text{2,3-диметил-}\ldots
    \]
    а не просто «диметил».

    \item \warning{Сравнивать строение, не проверив состав.}

    Различающиеся молекулярные формулы означают, что вещества не являются изомерами.

    \item \warning{Переставлять двойную связь без пересчёта числа атомов водорода.}

    После любого изменения структуры обязательно проверьте валентности и формулу вещества.
\end{enumerate}

% -------------------------------------------------

\lessonsection{13. Что нужно уметь после занятия}

Проверьте себя.

\begin{itemize}
    \item Я могу объяснить, что такое алкены.
    \item Я знаю общую формулу ациклических моноалкенов:
    \[
    \mathrm{C_nH_{2n}}.
    \]
    \item Я отличаю изомеры от разных изображений одной молекулы.
    \item Я различаю изомерию углеродного скелета и положения двойной связи.
    \item Я умею выбирать главную цепь, содержащую связь \(\mathrm{C=C}\).
    \item Я нумерую цепь так, чтобы двойная связь получила наименьший номер.
    \item Я узнаю метильный и этильный заместители.
    \item Я правильно использую приставки ди-, три-, тетра-.
    \item Я могу назвать несложный разветвлённый алкен.
    \item Я могу построить структурную формулу по названию.
    \item Я проверяю валентность углерода и молекулярную формулу.
\end{itemize}

% =================================================
% ТРЕНИРОВОЧНЫЙ БЛОК
% =================================================
\clearpage

\begin{center}
{\LARGE\bfseries\color{darkaccent}Тренировочный блок}

\vspace{0.25em}

{\color{linegray}\rule{0.82\linewidth}{0.8pt}}

\vspace{0.4em}

\textcolor{textgray}{Алкены: формулы, изомерия и номенклатура}
\end{center}

\sublesson{Средний уровень}

\begin{enumerate}[label=\textbf{\arabic*.}]

\item Запишите молекулярную формулу ациклического алкена, содержащего \(7\) атомов углерода.

\item Назовите вещество:
\[
\mathrm{CH_2=CH-CH_2-CH_2-CH_3}.
\]

\item Определите, являются ли изомерами вещества
\[
\mathrm{CH_2=CH-CH_2-CH_2-CH_3}
\]
и
\[
\mathrm{CH_3-CH=CH-CH_2-CH_3}.
\]
Если являются, укажите вид структурной изомерии.

\end{enumerate}

\sublesson{Выше среднего}

\begin{enumerate}[label=\textbf{\arabic*.},resume]

\item Назовите вещество:
\[
\mathrm{CH_2=C(CH_3)-CH_2-CH_3}.
\]

\item Постройте сокращённую структурную формулу
\[
\text{2-метилпентена-1}.
\]

\item Назовите вещество:
\[
\mathrm{CH_3-CH=C(CH_3)-CH_2-CH_3}.
\]

\item Постройте сокращённую структурную формулу
\[
\text{3-метилгексена-2}.
\]

\item Назовите вещество:
\[
\mathrm{CH_2=C(CH_3)-CH(CH_3)-CH_2-CH_3}.
\]

\item Для вещества
\[
\mathrm{CH_3-C(CH_3)=C(CH_3)-CH_2-CH_3}
\]
определите:
\begin{itemize}
    \item молекулярную формулу;
    \item положение двойной связи;
    \item полное название.
\end{itemize}

\end{enumerate}

\sublesson{Сложный уровень}

\begin{enumerate}[label=\textbf{\arabic*.},resume]

\item Рассмотрите структуру
\[
\mathrm{CH_2=C(CH_2CH_3)-CH_2-CH(CH_3)-CH_3}.
\]

\begin{enumerate}[label=\alph*)]
    \item Выберите главную цепь по правилам номенклатуры алкенов.
    \item Объясните, почему нельзя выбрать более длинную цепь, если она не содержит двойную связь.
    \item Определите положения заместителей.
    \item Назовите вещество.
    \item Запишите его молекулярную формулу.
\end{enumerate}

\end{enumerate}

% =================================================
% ОТВЕТЫ
% =================================================
\clearpage

\begin{center}
{\LARGE\bfseries\color{darkaccent}Ответы}

\vspace{0.25em}

{\color{linegray}\rule{0.82\linewidth}{0.8pt}}
\end{center}

\vspace{0.7em}

\renewcommand{\arraystretch}{1.45}

\begin{table}[ht]
\centering
\caption{Ответы к тренировочному блоку}
\begin{tabularx}{\linewidth}{@{}c X@{}}
\toprule
\textbf{№} & \textbf{Ответ} \\
\midrule

1 &
\(\mathrm{C_7H_{14}}\). \\

2 &
Пентен-1. \\

3 &
Да. Оба вещества имеют формулу \(\mathrm{C_5H_{10}}\).
Это изомеры положения двойной связи: пентен-1 и пентен-2. \\

4 &
2-метилбутен-1. \\

5 &
\(\mathrm{CH_2=C(CH_3)-CH_2-CH_2-CH_3}\). \\

6 &
3-метилпентен-2. \\

7 &
\(\mathrm{CH_3-CH=C(CH_3)-CH_2-CH_2-CH_3}\). \\

8 &
2,3-диметилпентен-1. \\

9 &
Молекулярная формула:
\(\mathrm{C_7H_{14}}\).

Двойная связь занимает положение \(2\).

Название:
2,3-диметилпентен-2. \\

10 &
Главная цепь должна содержать двойную связь и включает \(5\) атомов углерода.

При нумерации от конца, ближайшего к двойной связи, получаем пентен-1.
Заместители: этил при атоме \(2\) и метил при атоме \(4\).

Название:
4-метил-2-этилпентен-1.

Молекулярная формула:
\(\mathrm{C_8H_{16}}\). \\

\bottomrule
\end{tabularx}
\end{table}

\vfill

\begin{center}
{\small\color{textgray}
Главная стратегия: сначала двойная связь и главная цепь, затем нумерация,
после этого --- заместители и окончательная проверка структуры.}
\end{center}

\end{document}